% Part: first-order-logic
% Chapter: proof-systems
% Section: sequent-calculus

\documentclass[../../../include/open-logic-section]{subfiles}

\begin{document}

\iftag{FOL}
      {\olfileid[tr]{fol}{prf}{seq}}
      {\olfileid[tr]{pl}{prf}{seq}}

\olsection{Sekant Hesabı}

Birçok !!{derivation} sistemi !!{sentence}s içeren düzenlemelerle çalışırken,
sekant hesabı \emph{sekantlarla} çalışır. Sekant, şu biçimde bir ifadedir:
\[
!A_1, \dots, !A_m \Sequent !B_1, \dots, !B_n.
\]
Başka bir deyişle sekant, sekant simgesi~$\Sequent$ ile ayrılan iki
!!{sentence}s dizisinden oluşur. Dizilerden herhangi biri boş olabilir.
Bir !!^a{derivation}, sekant hesabında bir sekant ağacıdır. En üstteki sekantlar
özel bir biçimdedir (bunlara “başlangıç sekantları” veya “aksiyomlar” denir)
ve diğer her sekant, hemen üstündeki sekantlardan bir çıkarım kuralıyla elde
edilir. Çıkarım kuralları ya sekantlardaki !!{sentence}s üzerinde işlem yapar
(bunları sol veya sağ tarafta ekler, çıkarır ya da yeniden sıralar) ya da
kuralın sonucuna karmaşık bir !!{formula} getirir. Örneğin $\LeftR{\land}$
kuralı, $!A, \Gamma \Sequent \Delta$ ifadesinden $!A \land !B, \Gamma
\Sequent \Delta$ ifadesinin çıkarılmasına izin verir. $\RightR{\lif}$ kuralı
ise herhangi bir $\Gamma$, $\Delta$, $!A$ ve~$!B$ için $!A, \Gamma \Sequent
\Delta, !B$ ifadesinden $\Gamma \Sequent \Delta, !A \lif !B$ ifadesinin
çıkarılmasına izin verir. (Özellikle $\Gamma$ ve~$\Delta$ boş olabilir.)

Sekant hesabına dayalı $\Proves$ bağıntısı şöyle tanımlanır: $\Gamma_0$'daki
her $!A$ için $!A\in\Gamma$ olacak ve kökünde $\Gamma_0 \Sequent !A$ sekantı
bulunan bir !!{derivation} var olacak biçimde bir $\Gamma_0$ dizisi varsa,
$\Gamma \Proves !A$ olur. $\Sequent !A$ sekantı için !!a{derivation} varsa
$!A$, sekant hesabında bir teoremdir. Örneğin aşağıdaki !!a{derivation},
$\Proves (!A \land !B) \lif !A$ olduğunu gösterir:
\begin{prooftree}
  \Axiom$!A \fCenter !A$
  \RightLabel{\LeftR{\land}}
  \UnaryInf$!A \land !B \fCenter !A$
  \RightLabel{\RightR{\lif}}
  \UnaryInf$\fCenter (!A \land !B) \lif !A$
\end{prooftree}

$\Gamma_0 \Sequent$ sekantı için !!a{derivation} varsa (burada
$\Gamma_0$'daki her $!A$, $\Gamma$'nın elemanıdır ve sekantın sağ tarafı
boştur), $\Gamma$ kümesi sekant hesabında tutarsızdır. Tutarsız bir kümeden
\RightR{\Weakening} kuralıyla her bir !!{sentence} için onu !!{derive} mümkündür.

Sekant hesabı 1930'larda Gerhard Gentzen tarafından geliştirildi. Sistematik
ve simetrik tasarımı sayesinde, !!{derivation}s kuramı geliştirmek için çok
yararlı bir biçimciliktir. Sekant hesabında !!{derivation}s bulmak görece
kolaydır; ancak bunları okumak ve ispatlarla bağlantılarını görmek çoğu zaman
zordur. Yine de sekant hesabı, !!{derivation}s kullanan çok zarif bir
yaklaşım olduğunu kanıtlamıştır ve birçok mantığın bu tür bir !!{derivation}
sistemi vardır.

\end{document}
